% Shared exercise chunk: l4-binary-count-terminal
\begin{exercisechunk}
\item \textbf{Terminal cost at fixed KL.}
\label[exercise]{ex:binary-count-terminal}
\exerciseprofile{2}{1}{1}{\excalculation\quad\exinterpretation}{%
Show that models with equal sequence KL can have terminal costs differing by a
factor of order \(T\).}
Use the setting and Markov generators from \cref{ex:binary-count}.
Set \(\varepsilon=1-e^{-\kappa/T}\) for fixed \(\kappa>0\).
Compute the two terminal costs and the ratio
\(\E_{Q_{\mathrm{bad}}}[c_{\mathrm{fin}}]/
\E_{Q_{\mathrm{ind}}}[c_{\mathrm{fin}}]\). Explain why this ratio is of order \(T\) even though the terminal cost range is
fixed. Which model attains the terminal-cost upper endpoint?

\end{exercisechunk}
